\chapter{Partir de zéro}

Jusqu'ici, nous avons toujours pris un modèle existant pour lui greffer des
pièces. Ce chapitre montre l'autre voie : construire un modèle vide et lui
apprendre la langue elle-même.

\section{Pourquoi c'est rarement le bon choix}

Soyons francs avant de commencer.

\begin{center}
\begin{tabular}{p{3.6cm}p{4.8cm}p{5.2cm}}
\toprule
 & \textbf{Affiner (LoRA)} & \textbf{Partir de zéro} \\
\midrule
Point de départ & Un modèle qui sait déjà lire et écrire & Des nombres tirés au
                  hasard \\
Corpus & Quelques milliers d'exemples & Des milliards de mots \\
Durée & Heures à jours & Semaines à mois \\
Matériel & Une carte de 8~Go & Plusieurs cartes professionnelles \\
Résultat & Un modèle qui parle bien \emph{et} connaît votre domaine & Un modèle
           qui parle mal, sauf effort colossal \\
\bottomrule
\end{tabular}
\end{center}

\begin{nkpiege}
« Je veux créer ma propre IA de zéro » est un souhait naturel et presque
toujours une erreur de départ. Vous dépenseriez des mois à réapprendre à une
machine ce que la langue française est — travail déjà fait, et mieux fait, par
ceux dont c'est le métier.
\end{nkpiege}

\section{Pourquoi le faire quand même}

Trois raisons valables, et elles ne sont pas économiques :

\begin{enumerate}
  \item \textbf{Comprendre.} Voir un modèle passer du charabia à des mots
        français apprend plus que dix articles. C'est la meilleure raison.
  \item \textbf{Une langue peu dotée.} Si votre langue n'est présente dans aucun
        modèle existant, il n'y a rien à affiner. C'est le cas de nombreuses
        langues africaines.
  \item \textbf{Un problème étroit.} Pour une tâche très spécifique — reconnaître
        une forme, classer des messages courts — un petit modèle entraîné de
        zéro peut battre un grand modèle généraliste, et tourner mille fois plus
        vite.
\end{enumerate}

\begin{nkretenir}
Partez de zéro pour \emph{apprendre}, ou parce que rien n'existe dans votre
langue. Jamais pour « faire mieux » qu'un modèle existant sur une tâche
générale : le rapport de forces n'est pas en votre faveur.
\end{nkretenir}

\section{L'outil : NKGptTrain}

Le dépôt fournit un GPT complet, entraînable de zéro, avec son propre découpage
en tokens.

\begin{nkcode}[]{Compiler}
jenga build --target NKGptTrain --config Release
\end{nkcode}

Il se configure par un \textbf{fichier}, ce qui présente un avantage décisif sur
une longue série d'essais : le fichier se garde, se relit, se compare, se range
à côté du modèle produit. Six mois plus tard, vous saurez encore avec quels
réglages ce modèle a été obtenu.

\begin{nkcode}[]{mon-gpt.cfg}
# --- Corpus ---
NK_GPT_DIR=./mes-textes
NK_GPT_CHARS=1200000

# --- Taille du modèle ---
NK_GPT_T=256        # longueur de contexte, en tokens
NK_GPT_D=384        # largeur : nombres par token
NK_GPT_H=6          # têtes d'attention
NK_GPT_L=6          # couches
NK_GPT_B=16         # exemples traités ensemble
NK_GPT_MERGES=4000  # taille du vocabulaire

# --- Entraînement ---
NK_GPT_STEPS=20000
NK_GPT_LR=3e-4
NK_GPT_WARMUP=500
NK_GPT_SAVEEVERY=1000

# --- Contrôle ---
NK_GPT_VALFRAC=0.05
NK_GPT_VALEVERY=500

# --- Fichiers ---
NK_GPT_SAVE=mon-gpt.nkgp
\end{nkcode}

\begin{nkcode}[]{Lancer}
NK_GPT_CONFIG=mon-gpt.cfg NKGptTrain
\end{nkcode}

\section{Choisir la taille de son modèle}

Quatre nombres déterminent tout. Voici comment les lire.

\begin{center}
\begin{tabular}{lp{10.2cm}}
\toprule
\textbf{T} & \textbf{Contexte} : combien de tokens le modèle voit d'un coup. À
             256, il ne peut pas relier le début d'un long paragraphe à sa fin. \\
\textbf{D} & \textbf{Largeur} : combien de nombres décrivent un token. C'est la
             richesse de sa représentation. \\
\textbf{H} & \textbf{Têtes} : combien de façons différentes de relier les mots
             entre eux. \textbf{D doit être divisible par H.} \\
\textbf{L} & \textbf{Profondeur} : combien de couches. C'est ce qui permet le
             raisonnement en plusieurs temps. \\
\bottomrule
\end{tabular}
\end{center}

\begin{center}
\begin{tabular}{lrrrrp{4.4cm}}
\toprule
 & \textbf{T} & \textbf{D} & \textbf{H} & \textbf{L} & \textbf{Ce qu'on obtient} \\
\midrule
Jouet & 128 & 128 & 4 & 4 & Du français qui ressemble à du français, sans
                            sens. Quelques heures. \\
Petit & 256 & 384 & 6 & 6 & Des phrases correctes, un sens fragile.
                            \textbf{Le bon point de départ.} \\
Moyen & 512 & 768 & 12 & 12 & Des paragraphes tenables. Plusieurs jours. \\
\bottomrule
\end{tabular}
\end{center}

\begin{nkpiege}
La tentation est de viser grand tout de suite. Faites l'inverse : commencez par
le modèle « jouet ». S'il ne produit rien de lisible après quelques heures,
c'est votre corpus ou vos réglages qu'il faut corriger — et vous l'aurez appris
en heures plutôt qu'en semaines.
\end{nkpiege}

\section{Le corpus, ici, c'est du texte brut}

Pas de \texttt{Question:} ni de \texttt{Reponse:} : le modèle apprend à
\emph{continuer} du texte, pas à répondre. Donnez-lui des fichiers texte, il les
lira à la suite.

\begin{nkcode}[]{Organisation}
mes-textes/
  romans.txt
  articles.txt
  documentation.txt
\end{nkcode}

\textbf{Combien ?} Le réglage \texttt{NK\_GPT\_CHARS} fixe le budget de
caractères. Un million de caractères ($\approx$ 200\,000 mots, deux romans)
suffit pour un modèle jouet. En dessous, le modèle apprendra le texte par cœur
au lieu d'apprendre la langue.

\section{Reprendre, ici aussi}

Le principe du chapitre~5 vaut pour cet outil :

\begin{nkcode}[]{Reprendre un GPT}
NK_GPT_RESUME=1
NK_GPT_LOAD=mon-gpt.nkgp
\end{nkcode}

Les checkpoints portent l'extension \texttt{.nkgp} et sont écrits tous les
\texttt{NK\_GPT\_SAVEEVERY} pas.

\section{Voir le modèle apprendre}

C'est le moment qui justifie tout le chapitre. Générez du texte à intervalles
réguliers, et regardez :

\begin{center}
\begin{tabular}{lp{10.0cm}}
\toprule
\textbf{Pas} & \textbf{Ce que le modèle produit} \\
\midrule
0 & Des caractères au hasard. \\
500 & Des suites de lettres qui ont la \emph{forme} de mots, des espaces au bon
      endroit. \\
2\,000 & De vrais mots français, mais assemblés sans grammaire. \\
10\,000 & Des phrases grammaticalement correctes, au sens vacillant. \\
50\,000 & Des paragraphes cohérents sur quelques lignes. \\
\bottomrule
\end{tabular}
\end{center}

\begin{nkretenir}
Cette progression est la meilleure leçon d'apprentissage automatique qui soit.
Le modèle apprend la \emph{forme} avant le \emph{sens} — les lettres, puis les
mots, puis la grammaire, puis, très lentement, le sens. Aucun cours ne fait
comprendre cela aussi bien que de le voir.
\end{nkretenir}

\section{Ce que vous ne devez pas espérer}

Un modèle entraîné chez vous sur quelques millions de caractères produira des
phrases françaises. Il ne saura pas répondre à vos questions, ne connaîtra aucun
fait, et se contredira d'une phrase à l'autre.

Ce n'est pas un échec : c'est ce que produit cette quantité de données. Les
modèles auxquels vous êtes habitué ont vu des milliers de milliards de mots.

\begin{nkretenir}
Partir de zéro apprend \emph{comment} un modèle apprend. Pour un modèle
\emph{utile}, revenez au chapitre~4 : affiner un modèle existant reste, et de
loin, le meilleur usage de votre machine et de votre temps.
\end{nkretenir}

\begin{nkexo}
Entraînez un modèle jouet (T\,=\,128, D\,=\,128, H\,=\,4, L\,=\,4) sur un livre
du domaine public. Faites-lui produire dix lignes de texte au pas 100, au pas
1\,000, puis au pas 10\,000. Collez les trois résultats côte à côte : vous
tiendrez, en une page, une démonstration de ce qu'est l'apprentissage.
\end{nkexo}
